\section{OBJECTIVE AND REVIEW OF ALL VERSIONS}
\textbf{Erd\H{o}s \#313; independent attempt, 2026-09-23.}
Are there infinitely many primary pseudoperfect numbers $M$, satisfying $1/M+\sum_{p\mid M}1/p=1$? I read \texttt{313.tex} and \texttt{313\_v2.tex}, which repeat the same local congruence restrictions. I test a more flexible recursive move than adjoining the single prime $M+1$.
\section{WORK: EXACT TWO-PRIME EXTENSION CRITERION}
Let $M$ be primary pseudoperfect, and let $p,q$ be distinct primes not dividing $M$. Then $Mpq$ is primary pseudoperfect if and only if
\[
(p-M)(q-M)=M^2+1.
\]
Indeed, substitute $\sum_{r\mid M}1/r=1-1/M$ into the defining equation for $Mpq$. What remains is
\[
\frac1p+\frac1q+\frac1{Mpq}=\frac1M,
\]
which after multiplication by $Mpq$ is $M(p+q)+1=pq$, equivalent to the factored equation. Positivity in the reciprocal equation forces $p,q>M$. Thus all such extensions are obtained by factoring $M^2+1=de$ into positive integers and testing whether $M+d$ and $M+e$ are distinct primes. They automatically do not divide $M$.
For example $M=2$ gives $M^2+1=5=1\cdot5$, yielding $p=3$, $q=7$ and the known solution $42$. For $M=6$, $37=1\cdot37$ yields $7,43$ and $1806$. These are explicit checks of the criterion, not new primary pseudoperfect numbers.
\section{ADVERSARIAL VERIFICATION AND REMAINING GAP}
The two new primes must be distinct to preserve the squarefree prime-divisor formula. Merely satisfying the congruences in the supplied versions does not guarantee the reciprocal sum is one; the factored identity is necessary and sufficient for this particular extension. If $M^2+1$ is prime, the only factor pair is $(1,M^2+1)$, requiring both $M+1$ and $M^2+M+1$ prime. There is no argument that successful factor pairs continue indefinitely, and the method need not generate every primary pseudoperfect number. The infinitude problem remains.
\section{SOURCES CHECKED}
Both local versions listed above; indexed official entry \url{https://www.erdosproblems.com/313}, checked on 2026-09-23. The extension criterion is proved directly and no claim of literature novelty is made.
\section{FINAL}
\textbf{UNRESOLVED.} Two-prime extension is reduced to an exact finite divisor-and-primality test; unbounded successful iteration is unproved.
\section{COMPLETION ESTIMATE}
Conservative subjective progress toward infinitude.
\textbf{COMPLETION: 10\%}
